\chapter{2007年陕西卷（理）}

\section{单选题}

\begin{problem}
在复平面内，复数 \(z = \frac{1}{2 + \i}\) 对应的点位于（\quad）

\choices
    {第一象限}
    {第二象限}
    {第三象限}
    {第四象限}
\end{problem}

\begin{problem}
已知全集 \(U = \{1, 2, 3, 4, 5\}\)，集合 \(A = \{x \in \Z \mid |x - 3| < 2\}\)，则集合 \(\complement_U A\) 等于（\quad）

\choices
    {\( \{1, 2, 3, 4\} \)}
    {\( \{2, 3, 4\} \)}
    {\( \{1, 5\} \)}
    {\( \{5\} \)}
\end{problem}

\begin{problem}
抛物线 \( y = x^{2} \) 的准线方程是（\quad）

\choices
    {\(4y + 1 = 0\)}
    {\(4x + 1 = 0\)}
    {\(2y + 1 = 0\)}
    {\(2x + 1 = 0\)}
\end{problem}

\begin{problem}
已知 \(\sin \alpha = \frac{\sqrt{5}}{5}\)，则 \(\sin^4 \alpha - \cos^4 \alpha\) 的值为（\quad）

\choices
    {\(-\frac{1}{5}\)}
    {\(-\frac{3}{5}\)}
    {\(\frac{1}{5}\)}
    {\(\frac{3}{5}\)}
\end{problem}

\begin{problem}
各项均为正数的等比数列 \(\{a_{n}\}\) 的前 \(n\) 项和为 \(S_{n}\). 若 \(S_{n} = 2\)，\(S_{3n} = 14\)，则 \(S_{4n}\) 等于（\quad）

\choices
    {80}
    {30}
    {26}
    {16}
\end{problem}

\begin{problem}
一个正三棱锥的四个顶点都在半径为1的球面上，其中底面的三个顶点在该球的一个大圆上，则该正三棱锥的体积是（\quad）

\choices
    {\(\frac{3\sqrt{3}}{4}\)}
    {\( \frac{\sqrt{3}}{3} \)}
    {\(\frac{\sqrt{3}}{4}\)}
    {\(\frac{\sqrt{3}}{12}\)}
\end{problem}

\begin{problem}
已知双曲线 \(C: \frac{x^2}{a^2} - \frac{y^2}{b^2} = 1\)（\(a > 0\)，\(b > 0\)），以 \(C\) 的右焦点为圆心且与 \(C\) 的渐近线相切的圆的半径是（\quad）

\choices
    {\(\sqrt{ab}\)}
    {\(\sqrt{a^2 + b^2}\)}
    {a}
    {\(b\)}
\end{problem}

\begin{problem}
若函数 \( f(x) \) 的反函数 \( f^{-1}(x) \)，则函数 \( f(x - 1) \) 与 \( f^{-1}(x - 1) \) 的图象可能是（\quad）

\choices
    {\choicebitmap[width=0.15\paperwidth]{img/2007/shaanxi_science/q8_optA_fixed.png}}
    {\choicebitmap[width=0.15\paperwidth]{img/2007/shaanxi_science/q8_optB_fixed.png}}
    {\choicebitmap[width=0.15\paperwidth]{img/2007/shaanxi_science/q8_optC_fixed.png}}
    {\choicebitmap[width=0.15\paperwidth]{img/2007/shaanxi_science/q8_optD_fixed.png}}
\end{problem}

\begin{problem}
\begin{circlelist}
\item 四个非零实数 \(a\)，\(b\)，\(c\)，\(d\) 依次成等比数列的充要条件是 \(ad = bc\)；
\item 设 \(a\)，\(b \in \R\)，且 \(ab \neq 0\).
\item 若 \( f(x)=\log_{2}x \)，则 \( f(|x|) \) 是偶函数.
\end{circlelist}
若 \(\frac{a}{b} < 1\)，则 \(\frac{b}{a} > 1\)；
其中不正确的序号是（\quad）

\choices
    {\circled{1}\circled{2}\circled{3}}
    {\circled{1}\circled{2}}
    {\circled{2}\circled{3}}
    {\circled{1}\circled{3}}
\end{problem}

\begin{problem}
已知平面 \(\alpha\) \(\parallel\) 平面 \(\beta\)，直线 \(m \subset \alpha\)，直线 \(n \subset \beta\)，点 \(A \in m\)，点 \(B \in n\)，记点 \(A\)，\(B\) 之间的距离为 \(a\)，点 \(A\) 到直线 \(n\) 的距离为 \(b\)，直线 \(m\) 和 \(n\) 的距离为 \(c\)，则（\quad）

\choices
    {\(b \leqslant c \leqslant a\)}
    {\(a \leqslant c \leqslant b\)}
    {\(c \leqslant a \leqslant b\)}
    {\(c \leqslant b \leqslant a\)}
\end{problem}

\begin{problem}
\(f(x)\) 是定义在 \((0, + \infty)\) 上的非负可导函数，且满足 \(x f'(x) + f(x) \leqslant 0\). 对任意正数 \(a\)，\(b\)，若 \(a < b\)，则必有（\quad）

\choices
    {\(a f(b)\leqslant b f(a)\)}
    {\(b f(a)\leqslant a f(b)\)}
    {\(a f(a)\leqslant f(b)\)}
    {\(b f(b)\leqslant f(a)\)}
\end{problem}

\begin{problem}
设集合 \(S = \{A_0, A_1, A_2, A_3\}\)，在 \(S\) 上定义运算 \(\oplus\) 为：\(A_i \oplus A_j = A_k\)，其中 \(k\) 为 \(i + j\) 被 4 除的余数，\(i\)，\(j = 0\)，\(1\)，\(2\)，\(3\). 则满足关系式 \((x \oplus x) \oplus A_2 = A_0\) 的 \(x\)（\(x \in S\)）的个数为（\quad）

\choices
    {4}
    {3}
    {2}
    {1}
\end{problem}

\section{填空题}

\begin{problem}
\(\lim\limits_{x\to 1}\left(\frac{2x + 1}{x^2 + x - 2} -\frac{1}{x - 1}\right) =\fillinblank{}\).
\end{problem}

\begin{problem}
已知实数 \(x\)，\(y\) 满足条件 \(\begin{cases} x - 2y + 4 \geqslant 0, \\ 2x + y - 2 \geqslant 0, \\ 3x - y - 3 \leqslant 0, \end{cases}\) 则 \(z = x + 2y\) 的最大值为\fillinblank{}.
\end{problem}

\begin{problem}
如图，平面内有三个向量 \(\overrightarrow{OA}\)，\(\overrightarrow{OB}\)，\(\overrightarrow{OC}\)，其中 \(\overrightarrow{OA}\) 与 \(\overrightarrow{OB}\) 的夹角为 \(120^{\circ}\)，\(\overrightarrow{OA}\) 与 \(\overrightarrow{OC}\) 的夹角为 \(30^{\circ}\)，且 \(\left|\overrightarrow{OA}\right| = \left|\overrightarrow{OB}\right| = 1\)，\(\left|\overrightarrow{OC}\right| = 2\sqrt{3}\). 若 \(\overrightarrow{OC} = \lambda \overrightarrow{OA} + \mu \overrightarrow{OB}\)（\(\lambda, \mu \in \R\)），则 \(\lambda + \mu\) 的值为\fillinblank{}.
\bitmapfigure[max width=0.34\linewidth]{img/2007/shaanxi_science/q15_fig1.png}
\end{problem}

\begin{problem}
安排 3 名支教教师去 6 所学校任教，每校至多 2 人，则不同的分配方案共有\fillinblank{} 种. （用数字作答）
\end{problem}

\section{解答题}

\begin{problem}
设函数 \( f(x) = \bs{a} \cdot \bs{b} \)，其中向量 \( \bs{a} = (m, \cos 2x) \)，\( \bs{b} = (1 + \sin 2x, 1) \)，\( x \in \R \)，且函数 \( y = f(x) \) 的图象经过点 \( \left(\frac{\pi}{4}, 2\right) \).
\begin{enumerate}
\item 求实数 \(m\) 的值；
\item 求函数 \( f(x) \) 的最小值及此时 \(x\) 值的集合.
\end{enumerate}
\end{problem}

\begin{problem}
某项选拔共有三轮考核，每轮设有一个问题，能正确回答问题者进入下一轮考核，否则即被淘汰. 已知某选手能正确回答第一，二，三轮的问题的概率分别为 \(\frac{1}{4}\)，\(\frac{3}{5}\)，\(\frac{2}{5}\)，且各轮问题能否正确回答互不影响.
\begin{enumerate}
\item 求该选手被淘汰的概率；
\item 该选手在选拔中回答问题的个数记为 \( \xi \)，求随机变量 \( \xi \) 的分布列与数学期望.
\end{enumerate}
\end{problem}

\begin{problem}
如图，在底面为直角梯形的四棱锥 \(P - ABCD\) 中，\(AD // BC\)，\(\angle ABC = 90^\circ\)，\(PA \perp\) 平面 \(ABCD\)，\(PA = 4\)，\(AD = 2\)，\(AB = 2\sqrt{3}\)，\(BC = 6\).

\begin{enumerate}
\item 求证：\(BD \perp\) 平面 \(PAC\)；
\item 求二面角 \(A - PC - D\) 的大小.
\end{enumerate}

\bitmapfigure[max width=0.32\linewidth]{img/2007/shaanxi_science/q19_fig1.png}

\end{problem}

\begin{problem}
设函数 \( f(x) = \frac{\e^x}{x^2 + a x + a} \)，其中 \( a \) 为实数.
\begin{enumerate}
\item 若 \( f(x) \) 的定义域为 \(\R\)，求 \(a\) 的取值范围；
\item 当 \( f(x) \) 的定义域为 \(\R\) 时，求 \( f(x) \) 的单调减区间.
\end{enumerate}
\end{problem}

\begin{problem}
已知椭圆 \(C: \frac{x^2}{a^2} + \frac{y^2}{b^2} = 1\)（\(a > b > 0\)）的离心率为 \(\frac{\sqrt{6}}{3}\)，短轴一个端点到右焦点的距离为 \(\sqrt{3}\).
\begin{enumerate}
\item 求椭圆 \(C\) 的方程；
\item 设直线 \(l\) 与椭圆 \(C\) 交于 \(A\)，\(B\) 两点，坐标原点 \(O\) 到直线 \(l\) 的距离为 \( \frac{\sqrt{3}}{2} \)，求 \( \triangle AOB \) 面积的最大值.
\end{enumerate}
\end{problem}

\begin{problem}
已知各项全不为零的数列 \(\{a_{k}\}\) 的前 \(k\) 项和为 \(S_{k}\)，且 \(S_{k} = \frac{1}{2} a_{k} a_{k + 1}\)（\(k \in \mathbf{N}^{*}\)），其中 \(a_{1} = 1\).
\begin{enumerate}
\item 求数列 \( \{a_{k}\} \) 的通项公式；
\item 对任意给定的正整数 \(n\)（\(n \geqslant 2\)），数列 \(\{b_k\}\) 满足 \(\frac{b_{k+1}}{b_k} = \frac{k-n}{a_{k+1}}\) （\(k = 1\)，\(2\)，\(\cdots\)，\(n-1\)），\(b_1 = 1\). 求 \(b_1 + b_2 + \cdots + b_n\).
\end{enumerate}
\end{problem}
